\chapter{Reconciling Mean Bias and Spectral Anisotropy: A Unified Causal Theory of Low-Bit Attention Instability}
\label{chap:reconcile_bias_anisotropy}

\section{Problem Definition}

Transformer attention couples three mathematical objects: (i) high-dimensional representations produced by a forward pass, (ii) dot-product logits that define routing, and (iii) a nonlinear normalization (softmax) that produces a probability simplex. In full precision, these components are numerically stable and empirically effective. In low-bit training, however, the same pipeline can become unstable when the boundary between linear operators and nonlinear operators is forced into FP4: the outputs of FP4 linear layers (or their FP4-quantized surrogates) are fed directly into nonlinear operators such as softmax, and the resulting attention distributions become inaccurate, especially in the long tail.

The purpose of this chapter is to reconcile two seemingly competing explanations for low-bit instability that appear in prior drafts of this manuscript. One explanation emphasizes \emph{mean bias} as the dominant driver of extreme values; the other emphasizes \emph{spectral anisotropy} (spike--tail structure) and softmax amplification. Our updated understanding is that these explanations are not in conflict: they operate at different stages of the pipeline and dominate different failure modes. Mean bias is the first-order driver of elementwise outliers in activations and of extreme second-moment statistics, while spike--tail anisotropy and dot-product noise govern logit-gap distortion and softmax amplification after mean has been removed.

We formalize this as a causal chain of claims. Each claim is supported by (i) a theoretical analysis that isolates the mechanism and (ii) quantitative results from controlled experiments on real model tensors.

\section{Notation and Measurement Protocols}

We use $H$ for hidden dimension, $L$ for sequence length, $d$ for per-head dimension, and $N$ for the number of token vectors in a batch. For attention, we consider query, key, and value tensors
\[
Q,K,V \in \mathbb{R}^{B\times h \times L \times d},
\]
and logits
\[
Z = \frac{QK^\top}{\sqrt{d}} \in \mathbb{R}^{B\times h \times L \times L},
\qquad
A = \mathrm{softmax}(Z/\tau).
\]
For activation matrices $X\in\mathbb{R}^{N\times H}$, we use the token-mean vector
\[
\mu \triangleq \frac{1}{N}\mathbf{1}^\top X \in \mathbb{R}^{H},
\qquad
M \triangleq \mathbf{1}\mu^\top,
\qquad
X_c \triangleq X - M.
\]
We quantify whether a mean is structural or noise-level using the noise-normalized statistic
\[
R \;\triangleq\;
\frac{\|\mu\|_2}
{\sqrt{\mathbb{E}\|x_{\mathrm{noise}}\|_2^2 / N}}
\;\approx\;
\frac{\|\mu\|_2}{\sigma / \sqrt{N}},
\]
Here $\sigma^2$ denotes the coordinate-wise variance under a mean-zero noise model,
so that $\sigma / \sqrt{N}$ is the expected standard error of the mean.
Thus $R \approx 1$ corresponds to noise-level mean,
whereas $R \gg 1$ indicates a coherent structural bias.

For low-bit quantization, we use blockwise symmetric quantization with percentile-based scaling and stochastic rounding. The quantizer determines a step size $\Delta$ proportional to a high percentile (e.g., $p99.9$) or max of $|x|$ within a block; stochastic rounding keeps $\mathbb{E}[\varepsilon\mid x]=0$ while injecting variance of order $\Delta^2$.

For softmax error, we use KL divergence to the FP16/FP32 baseline attention matrix, along with tail log-error measured on low-probability entries. We also track collapse metrics (entropy and $k_{90}$), 
where $k_{90}$ denotes the number of tokens required to account for
$90\%$ of the attention probability mass.
These metrics distinguish true simplex collapse from mere tail distortion.

\section{A Causal Chain of Instability}

\subsection{Claim 1: A Coherent Mean Bias is Widespread in Forward Activations but Absent in Backward Gradient Outputs}

\paragraph{Mechanism.}
Forward activations are generated by a sequence of linear maps and nonlinear operators (e.g., GELU/SwiGLU, softmax attention). Many such operators are non-odd and therefore do not preserve zero-mean distributions. Even when inputs are centered, non-odd nonlinearities induce $\mathbb{E}[\phi(z)]\neq 0$ for symmetric $z$, creating a coherent drift in feature space. Residual pathways then accumulate these drifts across depth. In contrast, backward gradient outputs (error signals) are empirically mean-neutral; their mean is comparable to the finite-sample noise floor. This creates a forward/backward asymmetry: large coherent bias in forward representations but not in backward signals.

\paragraph{Quantitative evidence.}
On real tensors extracted from trained models, we measured noise-normalized mean magnitudes $R$. For key states in attention, using token-mean removal over the token dimension, we observed $R_{\text{raw}}\approx 127.9$ and $R_{\text{centered}}\approx 0$, confirming a strong coherent mean in $K$ that can be removed exactly by centering. For gradient-output tensors in earlier experiments, mean ratios remained near $R\approx 1$ (noise-level), indicating the absence of a persistent coherent mean. This aligns with our earlier activation-versus-gradient-output measurements: the forward path produces large coherent mean structure; the backward path does not.

\subsection{Claim 2: Mean Bias is the First-Order Driver of Elementwise Activation Outliers}

\paragraph{Mechanism.}
Elementwise outliers are governed by coordinatewise extremes, not by global energy. A coherent mean $M=\mathbf{1}\mu^\top$ shifts every coordinate of every token vector by $\mu_j$. In high dimension, even modest per-coordinate bias produces large coherent $\ell_2$ and large probability of exceeding extreme thresholds, because the bias is aligned and does not cancel. This is fundamentally different from spike variance: low-rank variance is coherent in a subspace but not necessarily dominant in coordinatewise extremes.

A minimal probabilistic model makes this explicit. Consider a scalar coordinate
\[
X_{ij} = \mu_j + Z_{ij},
\]
where $Z_{ij}$ is mean-zero sub-Gaussian with scale $\sigma$. For thresholds $t$ near $|\mu_j|$, the exceedance probability is dominated by the shift:
\[
\Pr(|X_{ij}|>t) = \Pr(|\mu_j + Z_{ij}|>t) \approx \Pr(Z_{ij}>t-\mu_j) + \Pr(Z_{ij}< -t-\mu_j).
\]
When $|\mu_j|\gg \sigma$, the exceedance event becomes controlled by the deterministic offset rather than rare fluctuations. In contrast, for mean-free $Z_{ij}$, $\Pr(|Z_{ij}|>t)$ decays as $\exp(-ct^2/\sigma^2)$.

\paragraph{Quantitative evidence.}
We decomposed activations into mean, spike, and tail components
\[
X = M + X_{\mathrm{spike}} + X_{\mathrm{tail}},
\]
where $X_{\mathrm{spike}}$ is the projection of $X_c$ onto the top-$k$ singular directions (typically $k\approx 1\%H$), and $X_{\mathrm{tail}}$ is the residual. For outliers defined as the top $0.1\%$ entries by $|X_{ij}|$, we computed per-outlier squared share contributions
\[
\rho^{(\mathrm{mean})}_{ij} = \frac{M_{ij}^2}{X_{ij}^2},\qquad
\rho^{(\mathrm{spike})}_{ij} = \frac{(X_{\mathrm{spike}})_{ij}^2}{X_{ij}^2},\qquad
\rho^{(\mathrm{tail})}_{ij} = \frac{(X_{\mathrm{tail}})_{ij}^2}{X_{ij}^2}.
\]
Because $X = M + X_{\mathrm{spike}} + X_{\mathrm{tail}}$,
the squared denominator $X_{ij}^2$ contains cross terms such as
$2 M_{ij}(X_{\mathrm{spike}})_{ij}$.
Therefore the $\rho$ fractions do not sum exactly to $1$.
However, for the top $0.1\%$ activation outliers,
$M_{ij}^2$ empirically dominates the magnitude,
rendering cross terms negligible in practice.

Across layers (e.g., layer 3 vs layer 21) and training stages (e.g., 10k vs 180k), we observed that among top outliers the mean share concentrates near one, with typical averages in the range $0.85$--$0.90$ of squared magnitude. Spike contributes a secondary fraction and tail is negligible. This establishes a robust empirical fact: \emph{activation outliers are mean-dominated}.

\subsection{Claim 3: After Mean Removal, Spike Does \emph{Not} Necessarily Dominate Coordinatewise Outliers, but It Governs Spectral Anisotropy and Tail SNR Under Quantization}

\paragraph{Mechanism.}
Subtracting $M$ removes the coherent rank-one shift, transforming outlier behavior from bias-driven to variance-driven. In $X_c$, coordinatewise extremes are no longer guaranteed to align with any single direction; they arise from a combination of structured spike variance and high-dimensional tail fluctuations. Consequently, it is incorrect to assert that ``the spike dominates outliers'' after mean removal in the same sense as mean bias. What remains true is that spike--tail anisotropy governs how quantization noise distributes relative to semantic energy, and thus governs tail signal-to-noise ratio (SQNR).

In the mean-free regime, consider a row vector $z\in\mathbb{R}^H$ decomposed into spike and tail components relative to a spike basis $U_k$:
\[
z = z_{\mathrm{sp}} + z_{\mathrm{tail}},
\qquad
z_{\mathrm{sp}} = U_kU_k^\top z,
\qquad
z_{\mathrm{tail}} = z - z_{\mathrm{sp}}.
\]
Blockwise quantization sets a step size $\Delta$ from a high percentile of $|z|$ within each block. Even if $z$ is mean-free, $\Delta$ can be controlled by whichever component produces large coordinatewise values within blocks. The key object that matters for long-tail learning is not whether spike dominates the maximum, but whether quantization noise projected into the tail overwhelms $z_{\mathrm{tail}}$:
\[
\mathrm{SQNR}_{\mathrm{tail}} \triangleq 10\log_{10}\frac{\|z_{\mathrm{tail}}\|_2^2}{\|\varepsilon_{\mathrm{tail}}\|_2^2},
\qquad
\varepsilon_{\mathrm{tail}} = z_{\mathrm{tail}} - \widehat{z}_{\mathrm{tail}}.
\]

\paragraph{Quantitative evidence.}
We carried out a spectral-domain analysis on mean-free attention logits $Z$ constructed from real $Q,K$ tensors. We formed logits
\[
Z = \frac{QK^\top}{\sqrt{d}},
\]
then row-centered them as $Z_c = Z - \mathrm{mean}_{\text{keys}}(Z)$, which is softmax-invariant. We computed a spike basis from the SVD of the row-stacked matrix of $Z_c$ and decomposed each row into spike and tail components. We then attributed the blockwise $p99.9$ scale $s$ to spike or tail by comparing $s_{\mathrm{spike}}$, $s_{\mathrm{tail}}$, and $s_{\mathrm{full}}$ across all blocks. In this logit space, the results showed that spike does \emph{not} dominate scale in general: the fraction of blocks with $s_{\mathrm{spike}} > s_{\mathrm{tail}}$ was $0.167$, and the median ratio $s_{\mathrm{spike}}/s_{\mathrm{full}}$ was $0.540$ while $s_{\mathrm{tail}}/s_{\mathrm{full}}$ was $0.794$. This demonstrates that, for attention logits, coordinatewise extremes under blockwise $p99.9$ scaling arise from a spike--tail competition rather than spike domination.

Despite this, spike whitening still improves tail SQNR. When we scaled the spike coefficients to match the average tail variance, the mean block scale decreased from $3.79$ to $2.97$, and the mean step size decreased from $0.541$ to $0.424$. Tail SQNR improved from $14.18$ dB to $16.37$ dB. This confirms that spike anisotropy contributes substantially to quantization noise injection into the tail, even when the tail can dominate blockwise scale.

\subsection{Claim 4: Softmax is Hypersensitive to Low-Bit Logit Perturbations Because It Operates on Logit Gaps, Not Absolute Scale}

\paragraph{Mechanism.}
Softmax depends on logit \emph{differences}. For any $i,j$,
\[
\frac{A_i}{A_j} = \exp\!\left(\frac{z_i-z_j}{\tau}\right).
\]
A perturbation $\widehat{z} = z + \varepsilon$ induces ratio distortion
\[
\frac{\widehat{A}_i}{\widehat{A}_j} = \frac{A_i}{A_j}\cdot \exp\!\left(\frac{\varepsilon_i-\varepsilon_j}{\tau}\right).
\]
Thus the relevant stability parameter is $\|\varepsilon\|/\tau$ compared to typical logit gaps $|z_i-z_j|$. When quantization coarsens logits such that many meaningful gaps fall below the step size, attention routing becomes unstable. This instability first appears in the tail because small probabilities correspond to log-probabilities that are sensitive to additive noise:
\[
\Delta \log A_i \approx \frac{1}{\tau}\left(\varepsilon_i - \sum_j A_j\varepsilon_j\right).
\]
Therefore tail log-error is a direct indicator of softmax amplification of quantization noise.

\paragraph{Quantitative evidence (bitwidth sensitivity).}
We performed a bitwidth sweep quantizing $Q,K$ with stochastic rounding and measuring softmax divergence against the FP16 baseline. For both biased and mean-free inputs, errors remained small at 12--8 bits and increased sharply at 6 bits, with a cliff at 4 bits. For the biased baseline, KL increased from $2.0\times 10^{-3}$ at 6 bits to $3.82\times 10^{-2}$ at 4 bits, and tail log-error increased from $6.38\times 10^{-2}$ to $2.62\times 10^{-1}$. For mean-free inputs, improvements were modest: at 4 bits, KL reduced from $3.82\times 10^{-2}$ to $3.34\times 10^{-2}$ and tail log-error from $2.62\times 10^{-1}$ to $2.50\times 10^{-1}$. This confirms that mean removal helps but does not remove the 4-bit cliff, consistent with the gap-sensitivity mechanism.

\subsection{Claim 5: Quantizing $Q,K$ Before the Dot Product Introduces Multiplicative Noise That Harms Logit-Gap Geometry More Than Quantizing Logits After the Dot Product}

\paragraph{Mechanism.}
Let $\hat Q = Q+\Delta Q$ and $\hat K=K+\Delta K$. Then
\[
\widehat{Z} = \frac{\hat Q\hat K^\top}{\sqrt{d}}
= Z + \frac{\Delta QK^\top + Q\Delta K^\top}{\sqrt{d}} + \frac{\Delta Q\Delta K^\top}{\sqrt{d}}.
\]
The leading perturbation terms are multiplicative: quantization noise is multiplied by full-precision $Q$ or $K$, yielding structured, heteroscedastic logit noise. In contrast, if one computes $Z$ in high precision and then quantizes logits, the perturbation is additive in logit space and does not include the cross terms above. The two strategies therefore generate qualitatively different error structure. In particular, the cross terms can distort top1--top2 gaps even if the argmax identity is sometimes preserved.

\paragraph{Quantitative evidence (logit-only vs QK-first).}
We compared two strategies at FP4 in logit space. Quantizing $Q,K$ first and then computing $\widehat{Z}$ yielded Frobenius relative error $1.56\times 10^{-1}$ and RMSE $3.06\times 10^{-1}$ against the FP32 baseline logits. Computing $Z$ first and then quantizing logits yielded lower Frobenius error $1.25\times 10^{-1}$ and lower RMSE $2.45\times 10^{-1}$. Most importantly, the relative distortion of the top1--top2 gap was dramatically worse for QK-first: the mean relative gap error was $2.58$ for QK-first versus $0.73$ for logit-only quantization. This provides direct evidence that QK-first quantization corrupts the margin geometry that governs softmax routing.

\paragraph{Quantitative evidence (end-to-end softmax pipelines).}
We then compared the two pipelines end-to-end under a shared softmax baseline. The QK-first pipeline yielded at 4 bits KL $3.80\times 10^{-2}$ and tail log-error $2.63\times 10^{-1}$. The logits-first pipeline with rowwise logit quantization yielded slightly worse KL $4.63\times 10^{-2}$ but substantially improved tail log-error $1.85\times 10^{-1}$. A variant that attempted to match clipping behavior via percentile sweep produced similar conclusions: tail log-error improved consistently, while KL did not uniformly improve. These results indicate that logits-first quantization better preserves long-tail attention mass, but logits themselves are heavy-tailed and require better logit-specific quantization schemes to reduce head distortion.

\subsection{Claim 6: Mean Bias, Spectral Anisotropy, Dot-Product Noise, and Softmax Amplification Form a Consistent Hierarchy of Failure Modes}

The preceding claims reconcile the two previously competing chapters. Mean bias is the dominant driver of activation outliers and extreme second-moment inflation, because it induces a coherent coordinatewise shift that dominates elementwise extremes. After mean removal, spike does not necessarily dominate coordinatewise outliers; instead, spike--tail anisotropy governs how quantization noise interacts with semantic energy, as reflected by tail SQNR. In attention, the most damaging low-bit effect is not mean bias per se, but logit-gap distortion induced by quantizing $Q,K$ before the dot product and by quantizing heavy-tailed logits. Softmax then amplifies these errors, particularly in the tail. These mechanisms are complementary and operate at distinct points in the computation graph.

\section{Practical Implications and Next Experiments}

Our up-to-date understanding implies that there is no single ``magic'' fix for FP4 attention. Mean removal is a first-order mitigation for activation outliers and optimizer second moments, and it modestly improves FP4 softmax robustness. However, the remaining failure modes are dominated by logit-gap preservation and by the difficulty of representing heavy-tailed logits and probability simplices in FP4. The most informative next experiments are therefore those that directly target these remaining mechanisms.

First, we should characterize the logit distribution and clipping dynamics more precisely by sweeping the percentile-based scale selection for logits quantization and measuring the tradeoff between KL, tail log-error, clipping fraction, and logit-gap distortion. Second, we should test logit soft-capping, which bounds $\|Z\|_\infty$ and directly reduces clipping and step size, and QK-normalization, which bounds $Q,K$ magnitude and reduces multiplicative cross-term noise. Third, we should evaluate probability-domain codecs (log-ratio encodings) when softmax outputs must be transmitted in low precision, since naive FP4 simplex quantization collapses the tail. Finally, we should propagate these attention-level improvements into end-to-end training, measuring not only softmax errors but also the downstream effect on value aggregation $AV$ and on gradient variance.

The reconciliation achieved in this chapter provides a coherent lens for interpreting all empirical results to date. Mean bias explains why outliers and second-moment extremes are so severe in activations, while spike--tail anisotropy and logit-gap sensitivity explain why FP4 attention remains difficult even after centering. Together, these mechanisms define the design space for truly stable low-bit transformers.